\secnumbersection{VALIDACIÓN DE LA SOLUCIÓN}

\textcolor{pink}{Se debe validar la solución propuesta. Esto significa probar o demostrar que la solución propuesta es válida para el entorno donde fue planteada.}

\textcolor{pink}{Tradicionalmente es una etapa crítica, pues debe comprobarse por algún medio que vuestra propuesta es básicamente válida. En el caso de un desarrollo de software es la construcción y sus pruebas; en el caso de propuestas de modelos, guías o metodologías podrían ser desde la aplicación a un caso real hasta encuestas o entrevistas con especialistas; en el caso de mejoras de procesos u optimizaciones, podría ser comparar la situación actual (previa a la memoria) con la situación final (cuando la memoria está ya implementada) en base a un conjunto cuantitativo de indicadores o criterios.}

\textcolor{pink}{\subsection{EJEMPLO DE COMO CITAR TABLAS}}

\textcolor{pink}{Se colocó una tabla que se puede referenciar también desde el texto (Ver tabla \ref{table:coloquios}).}

\begin{table}[h]
    \centering
    \caption{\label{table:coloquios} \textcolor{pink}{Coloquios del Ciclo de Charlas Informática.}} \textcolor{pink}{Fuente: Elaboración Propia.}
    \begin{tabular}{|p{7cm}|p{7cm}|}
        \hline
        \textcolor{pink}{Título Coloquio} & \textcolor{pink}{Presentador, País} \\
        \hline
        \textcolor{pink}{``Sensible, invisible, sometimes tolerant, heterogeneous, decentralized and interoperable... and we still need to assure its quality...'''} & \textcolor{pink}{Guilherme Horta Travassos, Brasil.}\\
        \hline
        \textcolor{pink}{``Dispersed Multiphase Flow Modeling: From Environmental to Industrial Applications'''} & \textcolor{pink}{Orlando Ayala, EE.UU.}\\
        \hline
        \textcolor{pink}{``Líneas de Producto Software Dinámicas para Sistemas atentos el Contexto'''} & \textcolor{pink}{Rafael Capilla, España.}\\
        \hline
        \textcolor{pink}{...} & \textcolor{pink}{...} \\
        \hline
    \end{tabular}
\end{table}

\subsection{ENFOQUE DE LA VALIDACIÓN}

A fin de validar el desarrollo, se realizaron pruebas con usuarios reales con discapacidad visual sobre el sistema ya construido. Se busca medir qué tan bien el Relator resuelve tareas reales de interpretación de imágenes en sitios web, integrando las consultas al Relator como parte del flujo de trabajo habitual del usuario con su lector de pantalla.

Dada la índole innovadora del trabajo, resulta difícil realizar una comparación cuantitativa formal —por ejemplo, contrastar los resultados de las tareas contra el lector de pantalla—, ya que los resultados esperados de uno y otro difieren demasiado entre sí. Por ello se opta por una prueba directa con el usuario, consultando su nivel de satisfacción tras completar cada tarea, llevando algo cualitativo (la sensación de satisfacción) a un valor cuantificable (un número entero).

\subsection{DISEÑO DEL PROCEDIMIENTO DE PRUEBA}

Se establece un número objetivo de entre 3 y 5 usuarios. Se exige que todos presenten discapacidad visual severa o ceguera total, que sean usuarios habituales de lectores de pantalla y que tengan un nivel de manejo computacional intermedio o avanzado.

Como entorno de pruebas se utiliza un subconjunto del corpus del benchmark (ver \ref{sec:benchmark}): al menos 2 sitios de categorías distintas y al menos 2 tareas por cada sitio. Además, se exige utilizar al menos una vez la modalidad por voz (VTV).

El detalle de los resultados por tarea se presenta en una tabla en el anexo. \textcolor{red}{REFERENCIA PENDIENTE: insertar \ref{...} a la tabla del anexo correspondiente.}

\subsection{MÉTRICA E INSTRUMENTO DE MEDICIÓN}

\subsubsection{Puntaje de satisfacción del usuario}

Como se mencionó, al terminar el flujo de trabajo de cada tarea se le pregunta al usuario su nivel de satisfacción. Este es un número entero entre 1 y 5, denominado puntaje de satisfacción (PS), donde 1 corresponde a insatisfacción total (el Relator alucinó información o su respuesta fue de nula ayuda para completar la tarea), 3 a satisfacción parcial (el Relator cumplió con lo mínimo solicitado, pero hay espacio de mejora) y 5 a satisfacción total (la respuesta fue la esperada y ayudó a completar la tarea sin problemas).

\subsubsection{Criterio de completitud}

Se definen como tareas parciales o completadas todas aquellas que obtengan un PS mayor o igual a 3. Este criterio se comenta a los usuarios al momento de consultar su puntaje, por lo que se asume una correlación directa entre el valor entregado y la completitud de la tarea.

\subsubsection{Estándar mínimo}

Se considera una sesión aprobada si, sobre el total de PS, se obtiene al menos un 55\% de tareas parciales o completadas (PS $\geq$ 3) y un promedio de al menos 55\%, equivalente a un PS promedio $\geq$ 2,75. El estándar mínimo por sesión se extrapola a los resultados globales (sumando todas las tareas de todas las sesiones y calculando sobre el total). Por lo tanto, para validar la solución propuesta se espera obtener al menos un 55\% de tareas aprobadas y un PS promedio global mayor a 2,75.

\textcolor{red}{figura: fórmula matemática para el cálculo de métricas.}


\subsection{EJECUCIÓN}

\subsubsection{Participantes}

La validación se realizó finalmente con 3 participantes, dentro del rango objetivo definido. Los tres presentan ceguera total y usan JAWS como lector de pantalla. Se diferencian entre sí en las condiciones asociadas a su discapacidad visual y en su nivel de manejo computacional (uno con nivel medio y dos con nivel avanzado).

\subsubsection{Sesiones}

Sobre estos 3 participantes se realizaron 4 sesiones en total, ya que P2 realizó dos sesiones. De ellas, 3 se hicieron en modalidad texto (TTV) y 1 en modalidad voz (VTV).

\textcolor{red}{tabla: detalle de sesiones (acá o anexo?).}

\subsection{RESULTADOS GLOBALES}

En total se evaluaron 19 tareas a lo largo de las 4 sesiones. El promedio global de satisfacción fue de 2,84 sobre 5 (aproximadamente 56\%), y 11 de las 19 tareas (57,9\%) obtuvieron un puntaje igual o superior a 3. Estos resultados cumplen ambos umbrales definidos como estándar mínimo en la sección anterior.

\textcolor{red}{tabla: resultados globales de satisfacción por tarea y sesión.}

\subsection{ANÁLISIS DE PATRONES}

A partir de la revisión de los registros de las sesiones se identificaron patrones en el desempeño del sistema. Por categoría de sitio, el mejor desempeño se observó en sitios de e-commerce y redes sociales, mientras que el peor desempeño se dio en sitios educacionales, particularmente en portales de gestión interna que requieren autenticación. Por tipo de tarea, el sistema tuvo buen desempeño en tareas descriptivas o de extracción de información --por ejemplo, ``qué es'', ``qué hay'' o ``describe''--, y un desempeño deficiente en tareas procedimentales o de navegación --por ejemplo, ``cómo hago'', ``dónde voy'' o solicitudes de pasos a seguir--.

Entre los modos de falla identificados se cuentan: un sesgo de perspectiva vidente en las respuestas del modelo, alucinaciones cuando el elemento consultado no está visible en la captura, descripciones superficiales, y fallas del mecanismo de scroll \& stitch (ver Capítulo 3, sección de extensión de captura) al capturar ventanas emergentes (no pop-ups, sino ventanas independientes usualmente para una sola función, como por ejemplo pagos online o visualizar documentos).

En la modalidad de voz se observó un comportamiento particular: se reportó una satisfacción más alta en la misma tarea ejecutada por texto incluso cuando esta no se completaba. Esto sugiere que la experiencia de interacción por voz tiene un valor intrínseco, independiente del éxito de la tarea.

\subsection{RETROALIMENTACIÓN DE LOS USUARIOS}

Los participantes aportaron las siguientes sugerencias de mejora durante las sesiones:

\begin{itemize}
    \item Contar con una versión móvil y compatible con aplicaciones, no solo con navegadores de escritorio.
    \item Que el chat mantuviera memoria de la conversación, para poder seguir preguntando sobre un mismo elemento sin repetir contexto.
    \item Poder acceder o navegar directamente a un enlace ya identificado por el sistema.
    \item Mejorar el feedback del sistema de entrada por voz, para que el usuario sepa si lo que dice se está capturando correctamente.
\end{itemize}

\subsection{CONCLUSIÓN DE LA VALIDACIÓN}

Considerando que se cumplen ambos umbrales del estándar mínimo definido —un 57,9\% de tareas con puntaje igual o superior a 3, por sobre el 55\% requerido, y un promedio global de 2,84, por sobre el mínimo de 2,75—, se declara la propuesta de solución VALIDADA.

Es necesario matizar esta conclusión: los resultados no son contundentes, sino que se ubican apenas por sobre el mínimo exigido. Los modos de falla identificados en la sección anterior y las sugerencias reportadas por los propios usuarios se abordan como trabajo futuro.
